% answers_short.tex —— 2015 年·山东卷·文科 · 简短答案速查
% 由 tools/make_short_answers.py 自 answers.tex 的 \gksolhead/\gkanswer 提取，
% 并与 2015/2015_数学真题_山东文_答案要点.md 逐题比对后生成，勿手改（改源 MD 后重跑）。

\begin{gknotice}
\textbf{使用说明}：本节只给最终答案便于快速核对，完整推导见解析版。答案为编写者独立推导，\textbf{非官方标准答案}。
\end{gknotice}

\gkpart{选择题}
\begin{center}\begin{tabular}{ccccccccccc}
\toprule
题号 & 1 & 2 & 3 & 4 & 5 & 6 & 7 & 8 & 9 & 10 \\
\midrule
答案 & C & A & C & B & D & B & A & C & B & D \\
\bottomrule
\end{tabular}\end{center}

\gkpart{填空题}
\begin{itemize}[leftmargin=2.6em]
  \item[\textbf{11.}] \(13\)
  \item[\textbf{12.}] \(7\)
  \item[\textbf{13.}] \(\dfrac32\)
  \item[\textbf{14.}] \(\sqrt2\)
  \item[\textbf{15.}] \(2+\sqrt3\)
\end{itemize}

\gkpart{解答题}
\begin{itemize}[leftmargin=2.6em]
  \item[\textbf{16.}] (1) \(\dfrac13\)．；(2) \(\dfrac2{15}\)．
  \item[\textbf{17.}] \(\sin A=\dfrac{2\sqrt2}{3}\)，\(c=1\)．
  \item[\textbf{18.}] (1) \(BD\parallel\) 平面 \(FGH\)．；(2) 平面 \(BCD\perp\) 平面 \(EGH\)．
  \item[\textbf{19.}] (1) \(a_n=2n-1\)．；(2) \(T_n=\dfrac{(12n-4)4^n+4}{9}\)．
  \item[\textbf{20.}] (1) \(a=1\)．；(2) 存在，\(k=1\)．；(3) 最大值为 \(\dfrac4{e^2}\)，在 \(x=2\) 处取得．
  \item[\textbf{21.}] (1) \(\dfrac{x^2}{4}+y^2=1\)．；(2) ① \(\dfrac{|OQ|}{|OP|}=2\)；② 最大面积为 \(6\sqrt3\)．
\end{itemize}
